\documentclass[a4paper,11pt]{article}

\usepackage[a4paper,margin=2.1cm]{geometry}
\usepackage{fontspec}
\usepackage{polyglossia}
\setmainlanguage{russian}
\setmainfont{Arial}
\setmonofont{Arial}

\usepackage{xcolor}
\usepackage{listings}
\usepackage{enumitem}
\usepackage{titlesec}
\usepackage{hyperref}
\usepackage{fancyhdr}

\definecolor{accent}{HTML}{2457A6}
\definecolor{codebg}{HTML}{F5F7FA}
\definecolor{darkgray}{HTML}{39424E}

\hypersetup{
  colorlinks=true,
  linkcolor=accent,
  urlcolor=accent,
  pdftitle={Инструкция по проведению эксперимента}
}

\lstdefinestyle{terminal}{
  backgroundcolor=\color{codebg},
  basicstyle=\ttfamily\small,
  columns=fullflexible,
  keepspaces=true,
  frame=single,
  rulecolor=\color{black!12},
  frameround=tttt,
  breaklines=true,
  showstringspaces=false
}

\titleformat{\section}
  {\Large\bfseries\color{accent}}
  {}{0pt}{}
\titleformat{\subsection}
  {\large\bfseries\color{darkgray}}
  {}{0pt}{}

\pagestyle{fancy}
\fancyhf{}
\setlength{\headheight}{14pt}
\lhead{PMLDL · LLM Classification Fine-tuning}
\rhead{Экспериментальный протокол}
\cfoot{\thepage}
\renewcommand{\headrulewidth}{0.4pt}

\newcommand{\step}[1]{\item \textbf{#1}}
\newcommand{\placeholder}[1]{\texttt{\textless #1\textgreater}}

\begin{document}

\begin{center}
  {\Huge\bfseries\color{accent} Инструкция по проведению эксперимента}\\[0.45em]
  {\large От чистого клона до воспроизводимого результата и Pull Request}
\end{center}

\section*{Подготовка (один раз на компьютере)}

\begin{lstlisting}[style=terminal]
git clone https://github.com/kujifined/PMLDL-llm-classification-finetuning.git
cd PMLDL-llm-classification-finetuning

python3 -m venv .venv
source .venv/bin/activate
python -m pip install -r requirements-baseline.lock
python -m pip install -e . --no-deps
python -m pip install clearml nbclient nbformat
clearml-init
\end{lstlisting}

\noindent\textit{Это устанавливает окружение и подключает ClearML. Ключи хранятся только локально.}

\section*{Проведение эксперимента}

\begin{enumerate}[leftmargin=*,itemsep=0.85em]
  \step{Создайте эксперимент и новую ветку. Это создаёт уникальный ID и конфигурацию.}
  \begin{lstlisting}[style=terminal]
make new-experiment
  \end{lstlisting}
  Ответьте на вопросы мастера и сохраните выведенный \texttt{ID}.

  \step{Измените код эксперимента. Так фиксируется проверяемая идея.}
  Откройте созданный файл\\
  \texttt{output/jupyter-notebook/}\placeholder{ID}\texttt{\textunderscore\textunderscore}\placeholder{название}\texttt{.ipynb}\\
  и измените только
  функцию \texttt{train\_and\_evaluate}. Конфигурацию, seed, split и метрики
  вручную не меняйте.

  \step{Запустите smoke-тест и полный эксперимент одной командой. Это проверяет код перед дорогим запуском.}
  \begin{lstlisting}[style=terminal]
make run-experiment EXPERIMENT=<ID>
  \end{lstlisting}
  Сначала выполняется быстрый smoke-тест. При ошибке команда остановится,
  сохранит traceback и укажет проблемную ячейку. После исправления запустите
  ту же команду повторно.

  \step{Проверьте результат. Это подтверждает, что метрики и логирование сохранились.}
  Метрики и metadata находятся в \texttt{results/runs/}, большие файлы --- в
  \texttt{artifacts/}, а live-метрики --- в ClearML-проекте
  \texttt{PMLDL LLM Classification Finetuning}.

  \step{Зафиксируйте и отправьте эксперимент. Команда проверяет run и обновляет leaderboard.}
  \begin{lstlisting}[style=terminal]
make submit-experiment EXPERIMENT=<ID>
  \end{lstlisting}
  Команда валидирует run, обновляет \texttt{results/leaderboard.csv}, создаёт
  commit, отправляет ветку на GitHub и печатает ссылку для Pull Request.

  \step{Создайте Pull Request.}
  Откройте напечатанную ссылку GitHub и дождитесь успешного прохождения CI.
\end{enumerate}

\section*{Минимальный контроль перед PR}

\begin{itemize}[leftmargin=*,itemsep=0.35em]
  \item full-run завершён со статусом \texttt{completed};
  \item метрики присутствуют в \texttt{results/runs/}\placeholder{run\_id}\texttt{/metrics.json};
  \item результат виден в ClearML и локальном leaderboard;
  \item в Git не добавлены данные Kaggle, токены, checkpoints и `.env`.
\end{itemize}

\vfill
\begin{center}
  \small\color{darkgray}
  После успешного CI эксперимент готов к сравнению с остальными.
\end{center}

\end{document}
